\documentclass[12pt,a4paper]{article}

%---------- Packages ----------
\usepackage[margin=2.5cm]{geometry}
\usepackage{amsmath, amssymb, amsthm}
\usepackage{booktabs}
\usepackage{array}
\usepackage{hyperref}
\usepackage{titlesec}
\usepackage{parskip}
\usepackage{fancyhdr}
\usepackage{xcolor}
\usepackage{tcolorbox}
\usepackage{mdframed}
\usepackage{enumitem}
\usepackage{multicol}
\usepackage{graphicx}
\usepackage{tabularx}
\usepackage{mathtools}
\usepackage{tikz}
\usetikzlibrary{arrows.meta, positioning}

\tcbuselibrary{skins, breakable, theorems}

%---------- Colours ----------
\definecolor{workedblue}{RGB}{25, 85, 160}
\definecolor{workedlight}{RGB}{230, 240, 255}
\definecolor{practicegreen}{RGB}{20, 120, 60}
\definecolor{practicelight}{RGB}{230, 250, 238}
\definecolor{headergray}{RGB}{55, 55, 65}
\definecolor{subtopicgold}{RGB}{180, 130, 0}
\definecolor{subtopiclight}{RGB}{255, 248, 220}
\definecolor{answerbox}{RGB}{245, 245, 245}
\definecolor{locolor}{RGB}{90, 40, 140}

%---------- tcolorbox styles ----------
\tcbset{
	worked/.style={
		colback=workedlight,
		colframe=workedblue,
		fonttitle=\bfseries\color{workedblue},
		title={Worked Example},
		breakable,
		before upper={\setlength{\parskip}{4pt}},
		top=6pt, bottom=6pt, left=8pt, right=8pt
	},
	practice/.style={
		colback=practicelight,
		colframe=practicegreen,
		fonttitle=\bfseries\color{practicegreen},
		title={Practice Question},
		breakable,
		before upper={\setlength{\parskip}{4pt}},
		top=6pt, bottom=6pt, left=8pt, right=8pt
	},
	subtopic/.style={
		colback=subtopiclight,
		colframe=subtopicgold,
		fonttitle=\bfseries\color{subtopicgold},
		breakable,
		top=4pt, bottom=4pt, left=8pt, right=8pt
	},
	answer/.style={
		colback=answerbox,
		colframe=gray!50,
		fonttitle=\bfseries\color{gray!70!black},
		title={Answer Space},
		breakable,
		top=4pt, bottom=80pt, left=8pt, right=8pt
	}
}

%---------- Custom commands ----------
\newcommand{\lo}[1]{\textcolor{locolor}{\textbf{[LO #1]}}}
\newcommand{\ac}[1]{\textcolor{locolor}{\textit{[AC #1]}}}
\newcommand{\blank}{\underline{\hspace{3cm}}}
\newcommand{\longblank}{\underline{\hspace{6cm}}}
\newcommand{\answerlines}{%
	\vspace{0.3cm}
	\hrule\vspace{0.5cm}
	\hrule\vspace{0.5cm}
	\hrule\vspace{0.5cm}
	\hrule\vspace{0.5cm}
	\hrule\vspace{0.5cm}
	\hrule
}
\newcommand{\step}[1]{\noindent\textbf{Step #1.}\;}
\newcommand{\R}{\mathbb{R}}
\newcommand{\mat}[1]{\begin{bmatrix}#1\end{bmatrix}}
\newcommand{\aug}[2]{\left[\begin{array}{c|c}#1 & #2\end{array}\right]}
\newcommand{\norm}[1]{\left\lVert #1 \right\rVert}

%---------- Section formatting ----------
\titleformat{\section}[block]
{\Large\bfseries\color{headergray}}
{\thesection}{1em}{}[{\color{headergray}\titlerule[1.2pt]}]

\titleformat{\subsection}[block]
{\normalsize\bfseries\color{subtopicgold}}
{\thesubsection}{0.8em}{}

%---------- Headers and footers ----------
\pagestyle{fancy}
\fancyhf{}
\rhead{\textcolor{gray}{Linear Algebra — Comprehensive Quiz}}
\cfoot{\textcolor{gray}{Page \thepage}}
\renewcommand{\headrulewidth}{0.4pt}

\hypersetup{
	colorlinks=true,
	linkcolor=workedblue,
	urlcolor=workedblue,
	pdftitle={ Comprehensive Linear Algebra Quiz},
}

%=============================================================
\begin{document}
	%=============================================================
	
	%---------- Title page ----------
	\begin{titlepage}
		\centering
		\vspace{2cm}
		{\Huge\bfseries\color{headergray} Linear Algebra with Applications\\[0.5em]}
		{\Large\color{subtopicgold} Comprehensive Quiz\\[2em]}
		{\large\color{gray} \quad|\quad All Learning Outcomes\\[0.5em]}
		
		\vspace{2.5cm}
		
		{\large
			\begin{tabular}{rl}
				\textbf{Student Name:}   & \underline{\hspace{8cm}} \\[1.2em]
				\textbf{Date:}           & \underline{\hspace{8cm}} \\[1.2em]
			\end{tabular}
		}
		
		\vspace{2cm}
		
		\begin{tcolorbox}[colback=gray!8, colframe=gray!50, width=0.9\textwidth,
			fonttitle=\bfseries, title={Instructions}]
			\begin{itemize}[nosep]
				\item Each section contains a \textbf{Worked Example} with a complete step-by-step
				solution, followed by a similar \textbf{Practice Question} to complete.
				\item Show all working clearly. Answers without supporting steps will not receive full credit.
				\item Calculators are permitted for numerical checks only.
				\item Unless stated otherwise, all entries are real numbers.
			\end{itemize}
			\vspace{0.5em}
			\begin{center}
				\begin{tabular}{lccc}
					\toprule
					\textbf{LO} & \textbf{Topic} & \textbf{Max Marks} & \textbf{Marks Obtained} \\
					\midrule
					1 & Systems, Matrices \& Determinants & 24 & \underline{\hspace{2cm}} \\
					2 & Vector Spaces \& Transformations  & 16 & \underline{\hspace{2cm}} \\
					3 & Orthogonality \& Scalar Products  &  8 & \underline{\hspace{2cm}} \\
					4 & Eigenvalues \& Eigenvectors       &  8 & \underline{\hspace{2cm}} \\
					5 & Numerical Methods \& Canonical Forms & 16 & \underline{\hspace{2cm}} \\
					\midrule
					\multicolumn{2}{l}{\textbf{Total}} & \textbf{72} & \underline{\hspace{2cm}} \\
					\bottomrule
				\end{tabular}
			\end{center}
		\end{tcolorbox}
	\end{titlepage}
	
	\tableofcontents
	\newpage
	
	%=============================================================
	\section{Learning Outcome 1: Systems of Linear Equations, Matrix Arithmetic, and Determinants}
	%=============================================================
	
	%-------------------------------------------------------------
	\subsection{AC 1.1 \& 1.2 — Solving Linear Systems and Row Echelon Form}
	%-------------------------------------------------------------
	
	\begin{tcolorbox}[subtopic]
		\textbf{Assessment Criteria 1.1 \& 1.2:} Analyse and solve linear equations using matrices;
		evaluate row echelon form (REF) and reduced row echelon form (RREF) and their role in solving linear systems.
	\end{tcolorbox}
	
	\begin{tcolorbox}[worked]
		\textbf{Problem.} Solve the following system using Gaussian elimination. State the solution and
		verify it.
		\[
		\begin{aligned}
			2x_1 + 4x_2 - 2x_3 &= 2 \\
			x_1 + 3x_2 +  x_3 &= 7 \\
			3x_1 + 5x_2 - 3x_3 &= 3
		\end{aligned}
		\]
		
		\step{1} Write the augmented matrix $[A|\mathbf{b}]$:
		\[
		\left[\begin{array}{ccc|c}
			2 & 4 & -2 & 2 \\
			1 & 3 &  1 & 7 \\
			3 & 5 & -3 & 3
		\end{array}\right]
		\]
		
		\step{2} $R_1 \leftarrow \tfrac{1}{2}R_1$ (make the leading entry 1):
		\[
		\left[\begin{array}{ccc|c}
			1 & 2 & -1 & 1 \\
			1 & 3 &  1 & 7 \\
			3 & 5 & -3 & 3
		\end{array}\right]
		\]
		
		\step{3} $R_2 \leftarrow R_2 - R_1$;\quad $R_3 \leftarrow R_3 - 3R_1$:
		\[
		\left[\begin{array}{ccc|c}
			1 & 2 & -1 & 1 \\
			0 & 1 &  2 & 6 \\
			0 & -1 &  0 & 0
		\end{array}\right]
		\]
		
		\step{4} $R_3 \leftarrow R_3 + R_2$ (eliminate below second pivot):
		\[
		\left[\begin{array}{ccc|c}
			1 & 2 & -1 & 1 \\
			0 & 1 &  2 & 6 \\
			0 & 0 &  2 & 6
		\end{array}\right]
		\quad \leftarrow \text{Row Echelon Form}
		\]
		
		\step{5} $R_3 \leftarrow \tfrac{1}{2}R_3$; then back-substitute:
		\[
		x_3 = 3 \qquad
		x_2 = 6 - 2(3) = 0 \qquad
		x_1 = 1 - 2(0) + (-1)(-3) = 1 - 0 + 3 \;\to\; x_1 = 1 + 3 = 4
		\]
		Wait — re-check step 5:
		$x_1 = 1 - 2x_2 + x_3 = 1 - 2(0) + 3 = 4$.
		
		\step{6} \textbf{Solution:} $\displaystyle(x_1, x_2, x_3) = (4,\; 0,\; 3)$.
		
		\textbf{Verification:}
		\begin{align*}
			2(4)+4(0)-2(3) &= 8-6 = 2\;\checkmark \\
			(4)+3(0)+(3)   &= 7\;\checkmark \\
			3(4)+5(0)-3(3) &= 12-9 = 3\;\checkmark
		\end{align*}
		
		\textbf{Note on REF:} A matrix is in row echelon form when (i) all zero rows are at the bottom,
		(ii) each leading entry (pivot) is to the right of the pivot in the row above, and
		(iii) all entries below a pivot are zero. The REF in Step 4 exhibits all three properties.
	\end{tcolorbox}
	
	\begin{tcolorbox}[practice]
		\textbf{Problem.} \lo{1} \ac{1.1, 1.2} \hfill \textit{(8 marks)}
		
		Solve the following system by reducing the augmented matrix to row echelon form, then use
		back-substitution. Verify your solution.
		\[
		\begin{aligned}
			3x_1 - 6x_2 + 3x_3 &= 9 \\
			2x_1 -  x_2 + 5x_3 &= 7 \\
			5x_1 - 8x_2 + 3x_3 &= 11
		\end{aligned}
		\]
		\begin{enumerate}[label=(\alph*)]
			\item Write the augmented matrix and reduce it to row echelon form, clearly showing each elementary row operation. \hfill\textit{(4)}
			\item State whether the system has a unique solution, infinitely many solutions, or no solution. Justify your answer from the REF. \hfill\textit{(2)}
			\item Determine the solution (or solution set) and verify at least one equation. \hfill\textit{(2)}
		\end{enumerate}
	\end{tcolorbox}
	
	\begin{tcolorbox}[answer]
	\end{tcolorbox}
	
	\newpage
	
	%-------------------------------------------------------------
	\subsection{AC 1.3 — Matrix Notation and Operations}
	%-------------------------------------------------------------
	
	\begin{tcolorbox}[subtopic]
		\textbf{Assessment Criterion 1.3:} Apply matrix notation and arithmetic operations,
		including addition, scalar multiplication, matrix multiplication, and the transpose.
	\end{tcolorbox}
	
	\begin{tcolorbox}[worked]
		\textbf{Problem.} Given
		\[
		A = \mat{1 & -2 & 3 \\ 0 & 4 & -1 \\ 2 & 1 & 0},\qquad
		B = \mat{2 & 1 \\ -1 & 3 \\ 4 & 0},\qquad
		\mathbf{v} = \mat{1 \\ -1 \\ 2},
		\]
		compute (a) $AB$,\quad (b) $A\mathbf{v}$,\quad (c) $B^T B$.
		
		\step{1} \textbf{Compute $AB$ ($3\times3$ times $3\times2$ $=$ $3\times2$):}
		\[
		AB = \mat{1(2)+(-2)(-1)+3(4) & 1(1)+(-2)(3)+3(0) \\
			0(2)+4(-1)+(-1)(4)  & 0(1)+4(3)+(-1)(0) \\
			2(2)+1(-1)+0(4)     & 2(1)+1(3)+0(0)}
		= \mat{16 & -5 \\ -8 & 12 \\ 3 & 5}
		\]
		
		\step{2} \textbf{Compute $A\mathbf{v}$ ($3\times3$ times $3\times1$):}
		\[
		A\mathbf{v} = \mat{1(1)+(-2)(-1)+3(2) \\ 0(1)+4(-1)+(-1)(2) \\ 2(1)+1(-1)+0(2)}
		= \mat{1+2+6 \\ 0-4-2 \\ 2-1} = \mat{9 \\ -6 \\ 1}
		\]
		
		\step{3} \textbf{Compute $B^T B$ ($2\times3$ times $3\times2$ $=$ $2\times2$):}
		\[
		B^T = \mat{2 & -1 & 4 \\ 1 & 3 & 0}, \qquad
		B^T B = \mat{2(2)+(-1)(-1)+4(4) & 2(1)+(-1)(3)+4(0) \\
			1(2)+3(-1)+0(4)     & 1(1)+3(3)+0(0)}
		= \mat{21 & -1 \\ -1 & 10}
		\]
		
		Note that $B^T B$ is \textbf{symmetric} (as expected for any $B^T B$).
	\end{tcolorbox}
	
	\begin{tcolorbox}[practice]
		\textbf{Problem.} \lo{1} \ac{1.3} \hfill \textit{(6 marks)}
		
		Given
		\[
		P = \mat{2 & 1 & -1 \\ 3 & 0 & 2 \\ -1 & 4 & 1},\qquad
		Q = \mat{1 & -2 \\ 2 & 0 \\ -1 & 3},\qquad
		\mathbf{u} = \mat{1 \\ 2 \\ -1},
		\]
		compute:
		\begin{enumerate}[label=(\alph*)]
			\item $PQ$ \hfill\textit{(2)}
			\item $P\mathbf{u}$ \hfill\textit{(2)}
			\item $Q^T Q$ and comment on any special property of the result. \hfill\textit{(2)}
		\end{enumerate}
	\end{tcolorbox}
	
	\begin{tcolorbox}[answer]
	\end{tcolorbox}
	
	\newpage
	
	%-------------------------------------------------------------
	\subsection{AC 1.4 — Elementary Matrices and Partitioned Matrices}
	%-------------------------------------------------------------
	
	\begin{tcolorbox}[subtopic]
		\textbf{Assessment Criterion 1.4:} Synthesise solutions involving elementary matrices and
		partitioned (block) matrices.
	\end{tcolorbox}
	
	\begin{tcolorbox}[worked]
		\textbf{Problem.}
		\begin{enumerate}[label=(\alph*), nosep]
			\item Write down the $3\times3$ elementary matrix $E$ that multiplies row 2 by $-3$ and adds it to row 3.
			\item Express the row operation applied to $A = \mat{2&1&0\\1&3&2\\4&0&1}$ as a matrix product $EA$, and compute $EA$.
			\item Find $E^{-1}$ and verify $E^{-1}E = I$.
		\end{enumerate}
		
		\step{1} \textbf{Elementary matrix $E$:} Start with $I_3$ and replace the $(3,2)$ entry with $-3$
		(add $-3 \times$ row 2 to row 3):
		\[
		E = \mat{1 & 0 & 0 \\ 0 & 1 & 0 \\ 0 & -3 & 1}
		\]
		
		\step{2} \textbf{Compute $EA$:}
		\[
		EA = \mat{1&0&0\\0&1&0\\0&-3&1}\mat{2&1&0\\1&3&2\\4&0&1}
		= \mat{2&1&0\\1&3&2\\4+(-3)(1)&0+(-3)(3)&1+(-3)(2)}
		= \mat{2&1&0\\1&3&2\\1&-9&-5}
		\]
		This is exactly the result of $R_3 \leftarrow R_3 - 3R_2$.
		
		\step{3} \textbf{Inverse $E^{-1}$:} Reverse the operation: add $+3\times$ row 2 to row 3:
		\[
		E^{-1} = \mat{1&0&0\\0&1&0\\0&3&1}
		\]
		\textbf{Verify:}
		\[
		E^{-1}E = \mat{1&0&0\\0&1&0\\0&3&1}\mat{1&0&0\\0&1&0\\0&-3&1}
		= \mat{1&0&0\\0&1&0\\0&3-3&-3+1} = \mat{1&0&0\\0&1&0\\0&0&1} = I\;\checkmark
		\]
	\end{tcolorbox}
	
	\begin{tcolorbox}[practice]
		\textbf{Problem.} \lo{1} \ac{1.4} \hfill \textit{(6 marks)}
		
		\begin{enumerate}[label=(\alph*)]
			\item Write the $3\times3$ elementary matrix $F$ corresponding to swapping rows 1 and 3 (a row-interchange operation). \hfill\textit{(1)}
			\item Let $B = \mat{5&2&1\\-1&3&0\\2&-4&6}$.
			Compute $FB$ and describe the effect on $B$. \hfill\textit{(2)}
			\item Find $F^{-1}$. What do you notice, and why does this make sense geometrically? \hfill\textit{(2)}
			\item A matrix $M$ is partitioned as
			$M = \left[\begin{smallmatrix} A & B \\ C & D \end{smallmatrix}\right]$
			where $A = \mat{1&2\\3&4}$, $B = \mat{0\\1}$, $C = \mat{-1&0}$, $D = [2]$.
			Write out $M$ in full and state its dimensions. \hfill\textit{(1)}
		\end{enumerate}
	\end{tcolorbox}
	
	\begin{tcolorbox}[answer]
	\end{tcolorbox}
	
	\newpage
	
	%-------------------------------------------------------------
	\subsection{AC 1.5 \& 1.6 — Determinants and Applications}
	%-------------------------------------------------------------
	
	\begin{tcolorbox}[subtopic]
		\textbf{Assessment Criteria 1.5 \& 1.6:} Calculate determinants and assess their properties;
		analyse the role of determinants in solving linear systems (Cramer's Rule) and other applications.
	\end{tcolorbox}
	
	\begin{tcolorbox}[worked]
		\textbf{Problem.} Let $A = \mat{3 & 1 & -2 \\ 0 & 4 & 1 \\ 2 & -1 & 3}$.
		
		\begin{enumerate}[label=(\alph*), nosep]
			\item Calculate $\det(A)$ by cofactor expansion along the first column.
			\item State two properties of determinants demonstrated in your calculation.
			\item Use Cramer's Rule to solve $A\mathbf{x} = \mat{1\\2\\0}$.
		\end{enumerate}
		
		\step{1} \textbf{Cofactor expansion along column 1:}
		\[
		\det(A) = 3\cdot M_{11} - 0\cdot M_{21} + 2\cdot M_{31}
		\]
		where $M_{ij}$ is $(-1)^{i+j}$ times the $(i,j)$ minor.
		\begin{align*}
			M_{11} &= (+1)\det\mat{4&1\\-1&3} = (12+1) = 13 \\
			M_{21} &= (-1)\det\mat{1&-2\\-1&3} = (-1)(3-2) = -1 \quad(\text{not needed}) \\
			M_{31} &= (+1)\det\mat{1&-2\\4&1} = (1+8) = 9
		\end{align*}
		\[
		\det(A) = 3(13) + 2(9) = 39 + 18 = \boxed{57}
		\]
		
		\step{2} \textbf{Properties illustrated:}
		\begin{itemize}[nosep]
			\item The row with a zero entry ($R_2$) contributes zero to the expansion, simplifying computation.
			\item $\det(A) \neq 0$, so $A$ is \textbf{invertible} and the system $A\mathbf{x}=\mathbf{b}$ has a unique solution.
		\end{itemize}
		
		\step{3} \textbf{Cramer's Rule:} $x_i = \dfrac{\det(A_i)}{\det(A)}$, where $A_i$ replaces column $i$ with $\mathbf{b}$.
		\[
		A_1 = \mat{1&1&-2\\2&4&1\\0&-1&3},\quad
		A_2 = \mat{3&1&-2\\0&2&1\\2&0&3},\quad
		A_3 = \mat{3&1&1\\0&4&2\\2&-1&0}
		\]
		\begin{align*}
			\det(A_1) &= 1(12+1)-1(6-0)+(-2)(-2-0) = 13-6+4 = 11 \\
			\det(A_2) &= 3(6-0)-1(0-2)+(-2)(0-4) = 18+2+8 = 28 \\
			\det(A_3) &= 3(0+2)-1(0-4)+1(0-8) = 6+4-8 = 2
		\end{align*}
		\[
		x_1 = \frac{11}{57},\quad x_2 = \frac{28}{57},\quad x_3 = \frac{2}{57}
		\]
	\end{tcolorbox}
	
	\begin{tcolorbox}[practice]
		\textbf{Problem.} \lo{1} \ac{1.5, 1.6} \hfill \textit{(8 marks)}
		
		Let $B = \mat{2 & -1 & 0 \\ 1 & 3 & -2 \\ 4 & 0 & 1}$.
		\begin{enumerate}[label=(\alph*)]
			\item Calculate $\det(B)$ using cofactor expansion. Choose your row or column carefully to minimise computation. \hfill\textit{(3)}
			\item Without further calculation, what does your answer tell you about the invertibility of $B$? \hfill\textit{(1)}
			\item Use Cramer's Rule to solve the system $B\mathbf{x} = \mat{1\\-1\\2}$.
			Show the construction of each modified matrix $B_i$ and the calculation of its determinant. \hfill\textit{(4)}
		\end{enumerate}
	\end{tcolorbox}
	
	\begin{tcolorbox}[answer]
	\end{tcolorbox}
	
	\newpage
	
	%=============================================================
	\section{Learning Outcome 2: Vector Spaces, Subspaces, and Linear Transformations}
	%=============================================================
	
	%-------------------------------------------------------------
	\subsection{AC 2.1 \& 2.2 — Vector Spaces, Linear Independence, Basis, and Dimension}
	%-------------------------------------------------------------
	
	\begin{tcolorbox}[subtopic]
		\textbf{Assessment Criteria 2.1 \& 2.2:} Define vector spaces and subspaces with examples;
		analyse linear independence, basis, and dimension.
	\end{tcolorbox}
	
	\begin{tcolorbox}[worked]
		\textbf{Problem.}
		\begin{enumerate}[label=(\alph*), nosep]
			\item Determine whether $S = \{\mathbf{v}_1, \mathbf{v}_2, \mathbf{v}_3\}$ is linearly independent, where
			$\mathbf{v}_1=(1,2,-1)$, $\mathbf{v}_2=(2,0,3)$, $\mathbf{v}_3=(0,4,-5)$.
			\item Find the dimension of the subspace $W = \text{Span}(S)$.
			\item Determine whether $W$ is a subspace of $\R^3$ and state two subspace axioms it satisfies.
		\end{enumerate}
		
		\step{1} \textbf{Test for linear independence:} set $c_1\mathbf{v}_1+c_2\mathbf{v}_2+c_3\mathbf{v}_3=\mathbf{0}$ and row-reduce:
		\[
		\mat{1&2&0\\2&0&4\\-1&3&-5} \xrightarrow{R_2-2R_1,\;R_3+R_1}
		\mat{1&2&0\\0&-4&4\\0&5&-5} \xrightarrow{R_3+\frac{5}{4}R_2}
		\mat{1&2&0\\0&-4&4\\0&0&0}
		\]
		
		The system has a free variable ($c_3$). Setting $c_3 = 1$: $c_2 = 1$, $c_1 = -2$.
		Hence $-2\mathbf{v}_1 + \mathbf{v}_2 + \mathbf{v}_3 = \mathbf{0}$, so $S$ is \textbf{linearly dependent}.
		
		\textbf{Verify:} $-2(1,2,-1)+(2,0,3)+(0,4,-5)=(-2+2+0,\;-4+0+4,\;2+3-5)=(0,0,0)\;\checkmark$
		
		\step{2} \textbf{Dimension of $W$:} The row-reduced matrix has \textbf{2 non-zero rows} (2 pivots), so
		$\dim(W) = 2$. A basis for $W$ is $\{\mathbf{v}_1, \mathbf{v}_2\}$ (the pivot columns).
		
		\step{3} \textbf{$W$ is a subspace of $\R^3$} because it is defined as a span:
		\begin{itemize}[nosep]
			\item \textbf{Closure under addition:} if $\mathbf{u}, \mathbf{w} \in W$ then $\mathbf{u}=\sum\alpha_i\mathbf{v}_i$ and $\mathbf{w}=\sum\beta_i\mathbf{v}_i$, so $\mathbf{u}+\mathbf{w}=\sum(\alpha_i+\beta_i)\mathbf{v}_i\in W$.
			\item \textbf{Closure under scalar multiplication:} $\lambda\mathbf{u}=\sum(\lambda\alpha_i)\mathbf{v}_i\in W$.
		\end{itemize}
	\end{tcolorbox}
	
	\begin{tcolorbox}[practice]
		\textbf{Problem.} \lo{2} \ac{2.1, 2.2} \hfill \textit{(8 marks)}
		
		Let $S = \{\mathbf{w}_1, \mathbf{w}_2, \mathbf{w}_3, \mathbf{w}_4\}$ where
		$\mathbf{w}_1=(1,0,2,1)$, $\mathbf{w}_2=(0,1,-1,2)$, $\mathbf{w}_3=(2,1,3,4)$, $\mathbf{w}_4=(1,-1,3,-1)$.
		
		\begin{enumerate}[label=(\alph*)]
			\item Row-reduce the matrix whose rows are $\mathbf{w}_1,\ldots,\mathbf{w}_4$. Show each step. \hfill\textit{(3)}
			\item Determine whether $S$ is linearly independent or dependent. If dependent, express one vector as a linear combination of the others. \hfill\textit{(2)}
			\item State the dimension of $W = \text{Span}(S)$ and give a basis for $W$. \hfill\textit{(2)}
			\item Verify that $\mathbf{0} \in W$, and use this to confirm $W$ is a subspace of $\R^4$. \hfill\textit{(1)}
		\end{enumerate}
	\end{tcolorbox}
	
	\begin{tcolorbox}[answer]
	\end{tcolorbox}
	
	\newpage
	
	%-------------------------------------------------------------
	\subsection{AC 2.3 \& 2.4 — Linear Transformations and Matrix Representations}
	%-------------------------------------------------------------
	
	\begin{tcolorbox}[subtopic]
		\textbf{Assessment Criteria 2.3 \& 2.4:} Define and apply linear transformations; analyse
		their matrix representations and the concept of similarity.
	\end{tcolorbox}
	
	\begin{tcolorbox}[worked]
		\textbf{Problem.} Define $T:\R^3\to\R^2$ by $T\!\begin{pmatrix}x_1\\x_2\\x_3\end{pmatrix}=\begin{pmatrix}2x_1-x_2+x_3\\x_1+3x_3\end{pmatrix}$.
		
		\begin{enumerate}[label=(\alph*), nosep]
			\item Show that $T$ is a linear transformation.
			\item Find the standard matrix $[T]$ of $T$.
			\item Find $\ker(T)$ and $\text{range}(T)$, and verify the Rank-Nullity Theorem.
		\end{enumerate}
		
		\step{1} \textbf{Linearity:} For $T$ to be linear, $T(\alpha\mathbf{u}+\beta\mathbf{v})=\alpha T(\mathbf{u})+\beta T(\mathbf{v})$.
		\[
		T(\alpha\mathbf{u}+\beta\mathbf{v})
		= \mat{2(\alpha u_1+\beta v_1)-(\alpha u_2+\beta v_2)+(\alpha u_3+\beta v_3)\\
			(\alpha u_1+\beta v_1)+3(\alpha u_3+\beta v_3)}
		= \alpha T(\mathbf{u})+\beta T(\mathbf{v}) \;\checkmark
		\]
		(Each component is a linear combination of the inputs — the transformation is linear.)
		
		\step{2} \textbf{Standard matrix:} Apply $T$ to each standard basis vector $\mathbf{e}_i$:
		\[
		T(\mathbf{e}_1)=\mat{2\\1},\quad T(\mathbf{e}_2)=\mat{-1\\0},\quad T(\mathbf{e}_3)=\mat{1\\3}
		\Rightarrow [T]=\mat{2&-1&1\\1&0&3}
		\]
		
		\step{3} \textbf{Kernel (null space):} Solve $[T]\mathbf{x}=\mathbf{0}$:
		\[
		\mat{2&-1&1\\1&0&3}\xrightarrow{R_1\leftrightarrow R_2}\mat{1&0&3\\2&-1&1}
		\xrightarrow{R_2-2R_1}\mat{1&0&3\\0&-1&-5}\Rightarrow
		x_2=5x_3,\;x_1=-3x_3
		\]
		$\ker(T)=\text{Span}\left\{\mat{-3\\5\\1}\right\}$, so $\text{nullity}(T)=1$.
		
		\textbf{Range:} Two pivot columns $\Rightarrow \text{rank}(T)=2$; range $=\R^2$.
		
		\textbf{Rank-Nullity:} $\text{rank}+\text{nullity}=2+1=3=\dim(\R^3)$ \;\checkmark
	\end{tcolorbox}
	
	\begin{tcolorbox}[practice]
		\textbf{Problem.} \lo{2} \ac{2.3, 2.4} \hfill \textit{(8 marks)}
		
		Define $T:\R^3\to\R^3$ by $T\!\begin{pmatrix}x_1\\x_2\\x_3\end{pmatrix}=\begin{pmatrix}x_1+2x_2\\3x_1-x_3\\x_2+4x_3\end{pmatrix}$.
		
		\begin{enumerate}[label=(\alph*)]
			\item Find the standard matrix $[T]$. \hfill\textit{(2)}
			\item Find the kernel and range of $T$. State $\text{rank}(T)$ and $\text{nullity}(T)$ and verify the Rank-Nullity Theorem. \hfill\textit{(4)}
			\item Two matrices $A$ and $B$ are \textbf{similar} if there exists an invertible matrix $P$ such that $B = P^{-1}AP$.
			Given $[T]$ above and $P=\mat{1&0&0\\0&1&1\\0&0&1}$, compute $P^{-1}[T]P$.
			What does similarity preserve? \hfill\textit{(2)}
		\end{enumerate}
	\end{tcolorbox}
	
	\begin{tcolorbox}[answer]
	\end{tcolorbox}
	
	\newpage
	
	%=============================================================
	\section{Learning Outcome 3: Orthogonality and Scalar Products}
	%=============================================================
	
	%-------------------------------------------------------------
	\subsection{AC 3.1 \& 3.2 — Scalar Products, Orthogonality, and Least Squares}
	%-------------------------------------------------------------
	
	\begin{tcolorbox}[subtopic]
		\textbf{Assessment Criteria 3.1 \& 3.2:} Analyse scalar products and orthogonality in $\R^n$;
		apply orthogonality to least squares problems and construct orthonormal sets.
	\end{tcolorbox}
	
	\begin{tcolorbox}[worked]
		\textbf{Problem.} Let $\mathbf{u}_1=(1,1,0)$, $\mathbf{u}_2=(1,-1,1)$, $\mathbf{b}=(3,1,2)$.
		
		\begin{enumerate}[label=(\alph*), nosep]
			\item Show that $\{\mathbf{u}_1,\mathbf{u}_2\}$ is an orthogonal set, then normalise to get an orthonormal set.
			\item Project $\mathbf{b}$ onto $W = \text{Span}\{\mathbf{u}_1,\mathbf{u}_2\}$.
			\item Find the component of $\mathbf{b}$ orthogonal to $W$ and verify your decomposition.
		\end{enumerate}
		
		\step{1} \textbf{Orthogonality check:}
		$\mathbf{u}_1\cdot\mathbf{u}_2 = (1)(1)+(1)(-1)+(0)(1)=0$ \checkmark
		
		\textbf{Normalise:}
		$\norm{\mathbf{u}_1}=\sqrt{1+1+0}=\sqrt{2}$,\;
		$\norm{\mathbf{u}_2}=\sqrt{1+1+1}=\sqrt{3}$
		\[
		\hat{\mathbf{u}}_1=\tfrac{1}{\sqrt{2}}(1,1,0),\quad
		\hat{\mathbf{u}}_2=\tfrac{1}{\sqrt{3}}(1,-1,1)
		\]
		Orthonormal set: $\{\hat{\mathbf{u}}_1,\hat{\mathbf{u}}_2\}$.
		
		\step{2} \textbf{Orthogonal projection of $\mathbf{b}$ onto $W$:}
		\[
		\hat{\mathbf{b}} = \frac{\mathbf{b}\cdot\mathbf{u}_1}{\mathbf{u}_1\cdot\mathbf{u}_1}\mathbf{u}_1
		+ \frac{\mathbf{b}\cdot\mathbf{u}_2}{\mathbf{u}_2\cdot\mathbf{u}_2}\mathbf{u}_2
		\]
		\[
		\mathbf{b}\cdot\mathbf{u}_1=3+1+0=4,\quad \mathbf{u}_1\cdot\mathbf{u}_1=2
		\Rightarrow \text{coeff}_1=2
		\]
		\[
		\mathbf{b}\cdot\mathbf{u}_2=3-1+2=4,\quad \mathbf{u}_2\cdot\mathbf{u}_2=3
		\Rightarrow \text{coeff}_2=\tfrac{4}{3}
		\]
		\[
		\hat{\mathbf{b}} = 2(1,1,0)+\tfrac{4}{3}(1,-1,1)
		= (2,2,0)+(\tfrac{4}{3},-\tfrac{4}{3},\tfrac{4}{3})
		= \left(\tfrac{10}{3},\;\tfrac{2}{3},\;\tfrac{4}{3}\right)
		\]
		
		\step{3} \textbf{Orthogonal component:}
		$\mathbf{b}-\hat{\mathbf{b}}=(3-\tfrac{10}{3},\;1-\tfrac{2}{3},\;2-\tfrac{4}{3})=(-\tfrac{1}{3},\;\tfrac{1}{3},\;\tfrac{2}{3})$.
		
		\textbf{Verify:} $(-\tfrac{1}{3},\tfrac{1}{3},\tfrac{2}{3})\cdot(1,1,0)=-\tfrac{1}{3}+\tfrac{1}{3}=0\;\checkmark$;\quad
		$(-\tfrac{1}{3},\tfrac{1}{3},\tfrac{2}{3})\cdot(1,-1,1)=-\tfrac{1}{3}-\tfrac{1}{3}+\tfrac{2}{3}=0\;\checkmark$
	\end{tcolorbox}
	
	\begin{tcolorbox}[practice]
		\textbf{Problem.} \lo{3} \ac{3.1, 3.2} \hfill \textit{(8 marks)}
		
		Let $\mathbf{v}_1=(2,1,-1)$, $\mathbf{v}_2=(1,-1,0)$, and $\mathbf{c}=(4,2,1)$.
		
		\begin{enumerate}[label=(\alph*)]
			\item Verify that $\{\mathbf{v}_1,\mathbf{v}_2\}$ is an orthogonal set, then construct an orthonormal set $\{\hat{\mathbf{v}}_1,\hat{\mathbf{v}}_2\}$. \hfill\textit{(2)}
			\item Compute the orthogonal projection $\hat{\mathbf{c}}$ of $\mathbf{c}$ onto $W=\text{Span}\{\mathbf{v}_1,\mathbf{v}_2\}$. \hfill\textit{(3)}
			\item Find the component $\mathbf{c}-\hat{\mathbf{c}}$ orthogonal to $W$ and verify it is perpendicular to both $\mathbf{v}_1$ and $\mathbf{v}_2$. \hfill\textit{(2)}
			\item Interpret the decomposition $\mathbf{c}=\hat{\mathbf{c}}+(\mathbf{c}-\hat{\mathbf{c}})$ in the context of a least-squares problem: if $A=\begin{bmatrix}\mathbf{v}_1 & \mathbf{v}_2\end{bmatrix}$ (columns), what does $\hat{\mathbf{c}}$ represent? \hfill\textit{(1)}
		\end{enumerate}
	\end{tcolorbox}
	
	\begin{tcolorbox}[answer]
	\end{tcolorbox}
	
	\newpage
	
	%=============================================================
	\section{Learning Outcome 4: Eigenvalues and Eigenvectors}
	%=============================================================
	
	%-------------------------------------------------------------
	\subsection{AC 4.1 \& 4.2 — Eigenanalysis and Applications}
	%-------------------------------------------------------------
	
	\begin{tcolorbox}[subtopic]
		\textbf{Assessment Criteria 4.1 \& 4.2:} Analyse eigenvalues and eigenvectors including
		complex eigenvalues; apply eigenvalues to solving linear systems and understanding transformations.
	\end{tcolorbox}
	
	\begin{tcolorbox}[worked]
		\textbf{Problem.} Let $A=\mat{4&-2\\1&1}$.
		
		\begin{enumerate}[label=(\alph*), nosep]
			\item Find the eigenvalues of $A$.
			\item Find an eigenvector for each eigenvalue.
			\item Diagonalise $A$: express $A = PDP^{-1}$ where $D$ is diagonal.
			\item Use the diagonalisation to compute $A^3$.
		\end{enumerate}
		
		\step{1} \textbf{Characteristic equation:} $\det(A-\lambda I)=0$:
		\[
		\det\mat{4-\lambda&-2\\1&1-\lambda}=(4-\lambda)(1-\lambda)+2=\lambda^2-5\lambda+6=(\lambda-2)(\lambda-3)=0
		\]
		Eigenvalues: $\lambda_1=2$, $\lambda_2=3$.
		
		\step{2} \textbf{Eigenvectors:}
		
		For $\lambda_1=2$: $(A-2I)\mathbf{x}=\mathbf{0}$:
		$\mat{2&-2\\1&-1}\mathbf{x}=\mathbf{0}\Rightarrow x_1=x_2$; eigenvector $\mathbf{p}_1=\mat{1\\1}$.
		
		For $\lambda_2=3$: $(A-3I)\mathbf{x}=\mathbf{0}$:
		$\mat{1&-2\\1&-2}\mathbf{x}=\mathbf{0}\Rightarrow x_1=2x_2$; eigenvector $\mathbf{p}_2=\mat{2\\1}$.
		
		\step{3} \textbf{Diagonalisation:}
		\[
		P=\mat{1&2\\1&1},\quad D=\mat{2&0\\0&3},\quad P^{-1}=\frac{1}{\det P}\mat{1&-2\\-1&1}=\mat{-1&2\\1&-1}
		\]
		\textbf{Verify:} $AP = PD$ iff $A\mathbf{p}_i=\lambda_i\mathbf{p}_i$:
		$A\mat{1\\1}=\mat{2\\2}=2\mat{1\\1}\;\checkmark$;\;
		$A\mat{2\\1}=\mat{6\\3}=3\mat{2\\1}\;\checkmark$
		
		\step{4} \textbf{$A^3 = PD^3P^{-1}$:}
		\[
		D^3=\mat{8&0\\0&27},\quad
		A^3=\mat{1&2\\1&1}\mat{8&0\\0&27}\mat{-1&2\\1&-1}
		=\mat{8&54\\8&27}\mat{-1&2\\1&-1}
		=\mat{46&-38\\19&-11}
		\]
	\end{tcolorbox}
	
	\begin{tcolorbox}[practice]
		\textbf{Problem.} \lo{4} \ac{4.1, 4.2} \hfill \textit{(8 marks)}
		
		Let $B=\mat{3&-5\\1&-1}$.
		\begin{enumerate}[label=(\alph*)]
			\item Find the characteristic polynomial of $B$ and compute its eigenvalues. Note: the discriminant is negative — state the complex eigenvalues in the form $a\pm bi$. \hfill\textit{(2)}
			\item For each eigenvalue, find a corresponding (complex) eigenvector. \hfill\textit{(2)}
			\item Describe geometrically what a transformation with complex eigenvalues $a\pm bi$ does to vectors in $\R^2$. (Hint: consider the modulus $r=\sqrt{a^2+b^2}$ and argument $\theta=\arctan(b/a)$.) \hfill\textit{(2)}
			\item Find $\text{tr}(B)$ and $\det(B)$ and relate them to the sum and product of the eigenvalues respectively. \hfill\textit{(2)}
		\end{enumerate}
	\end{tcolorbox}
	
	\begin{tcolorbox}[answer]
	\end{tcolorbox}
	
	\newpage
	
	%=============================================================
	\section{Learning Outcome 5: Numerical Methods and Canonical Forms}
	%=============================================================
	
	%-------------------------------------------------------------
	\subsection{AC 5.1 \& 5.2 — Floating-Point Numbers and Numerical Gaussian Elimination}
	%-------------------------------------------------------------
	
	\begin{tcolorbox}[subtopic]
		\textbf{Assessment Criteria 5.1 \& 5.2:} Analyse floating-point representation; apply
		Gaussian elimination with pivoting strategies and matrix norms.
	\end{tcolorbox}
	
	\begin{tcolorbox}[worked]
		\textbf{Problem.}
		\begin{enumerate}[label=(\alph*), nosep]
			\item Represent $x = 0.1$ in IEEE 754 single-precision floating-point (32-bit).
			Why can it not be stored exactly?
			\item Use Gaussian elimination \textbf{with partial pivoting} to solve:
			$\mat{0.001 & 1 \\ 1 & 2}\mat{x_1\\x_2}=\mat{1\\3}$.
			Contrast with naive elimination (no pivoting) to show the benefit.
		\end{enumerate}
		
		\step{1} \textbf{Floating-point representation of $0.1$:}
		$0.1 = 0.0\overline{0011}_2$ (the binary expansion is periodic/non-terminating).
		With 23 explicit mantissa bits, the stored value is approximately
		$0.10000000149\ldots \neq 0.1$, introducing round-off error.
		IEEE 754 single: sign $0$, exponent $-4+127=123=(01111011)_2$, mantissa truncated to 23 bits.
		
		\step{2} \textbf{Naive elimination (without pivoting):}
		Multiplier $m = 1/0.001 = 1000$. After elimination:
		$\mat{0.001 & 1 \\ 0 & 2-1000}=\mat{0.001 & 1 \\ 0 & -998}\mat{\ |\ \!1\\\ |\ -997}$
		
		$x_2 = -997/-998 \approx 0.999$;\; $x_1 = (1-1\cdot0.999)/0.001 \approx 1.0$.
		This is numerically unstable if 0.001 suffers catastrophic cancellation.
		
		\step{3} \textbf{With partial pivoting:} Swap rows (pivot on largest in column 1):
		\[
		\mat{1&2\\0.001&1}\mat{|\\|}\ \mat{3\\1}
		\xrightarrow{R_2-0.001R_1}
		\mat{1&2\\0&0.998}\mat{|\\|}\ \mat{3\\0.997}
		\]
		$x_2 = 0.997/0.998 \approx 0.999$;\; $x_1 = 3-2(0.999) = 1.002$.
		
		\textbf{Exact solution:} $A^{-1}$ gives $x_1\approx1.0$, $x_2\approx1.0$. Partial pivoting gives a more accurate result by avoiding division by a near-zero pivot.
		
		\textbf{$\infty$-norm of $A$:} $\norm{A}_\infty = \max(\text{row sums}) = \max(|0.001|+|1|,\;|1|+|2|) = 3$.
	\end{tcolorbox}
	
	\begin{tcolorbox}[practice]
		\textbf{Problem.} \lo{5} \ac{5.1, 5.2} \hfill \textit{(8 marks)}
		
		\begin{enumerate}[label=(\alph*)]
			\item A floating-point system uses 4-bit mantissa (including sign) and 3-bit exponent in base 2.
			Represent $x = 5.5$ in this system and identify the machine epsilon $\varepsilon_{\text{mach}}$.
			What is the smallest positive representable number in this system? \hfill\textit{(3)}
			\item Apply Gaussian elimination \textbf{with partial pivoting} to solve:
			$\mat{2&1&-1\\4&5&2\\-2&1&4}\mat{x_1\\x_2\\x_3}=\mat{8\\22\\-2}$.
			Show each pivot selection and row swap. Verify your solution. \hfill\textit{(5)}
		\end{enumerate}
	\end{tcolorbox}
	
	\begin{tcolorbox}[answer]
	\end{tcolorbox}
	
	\newpage
	
	%-------------------------------------------------------------
	\subsection{AC 5.3 \& 5.4 — Jordan Canonical Form and Nilpotent Operators}
	%-------------------------------------------------------------
	
	\begin{tcolorbox}[subtopic]
		\textbf{Assessment Criteria 5.3 \& 5.4:} Analyse nilpotent operators and the Jordan canonical
		form; apply canonical form concepts in context.
	\end{tcolorbox}
	
	\begin{tcolorbox}[worked]
		\textbf{Problem.} Let $N=\mat{0&1&0\\0&0&1\\0&0&0}$ and $A=\mat{3&1&0\\0&3&1\\0&0&3}$.
		
		\begin{enumerate}[label=(\alph*), nosep]
			\item Show that $N$ is nilpotent and find its index of nilpotency.
			\item Write $A$ in Jordan canonical form. Identify the Jordan block(s).
			\item Compute $e^A = e^{3I+N}$ using the matrix exponential series. (Hint: $e^{3I+N} = e^{3I}e^N$.)
		\end{enumerate}
		
		\step{1} \textbf{Nilpotency of $N$:}
		\[
		N^2 = \mat{0&0&1\\0&0&0\\0&0&0},\qquad N^3 = \mat{0&0&0\\0&0&0\\0&0&0} = O
		\]
		$N^2 \neq O$ but $N^3 = O$, so $N$ is nilpotent with \textbf{index} $k=3$.
		
		\step{2} \textbf{Jordan form of $A$:}
		The only eigenvalue of $A$ is $\lambda = 3$ (multiplicity 3).
		Since $A - 3I = N \neq 0$ and $(A-3I)^2 = N^2 \neq 0$ but $(A-3I)^3=0$,
		the Jordan block has size 3. The Jordan canonical form is:
		\[
		J = \mat{3&1&0\\0&3&1\\0&0&3}
		\]
		(Here $A$ is already in Jordan form; the single $3\times3$ block corresponds to eigenvalue $\lambda=3$.)
		
		\step{3} \textbf{Matrix exponential $e^A$:}
		Since $3I$ and $N$ commute, $e^{3I+N}=e^{3I}e^N$.
		\[
		e^{3I} = e^3 I,\qquad
		e^N = I + N + \tfrac{1}{2!}N^2 = \mat{1&1&\frac{1}{2}\\0&1&1\\0&0&1}
		\]
		\[
		e^A = e^3 \mat{1&1&\frac{1}{2}\\0&1&1\\0&0&1}
		\]
		
		\textbf{Significance:} The Jordan form reveals that $A$ has a defective eigenvalue (geometric multiplicity 1 $<$ algebraic multiplicity 3), meaning $A$ is \emph{not} diagonalisable over $\R$.
	\end{tcolorbox}
	
	\begin{tcolorbox}[practice]
		\textbf{Problem.} \lo{5} \ac{5.3, 5.4} \hfill \textit{(8 marks)}
		
		Let $C = \mat{2&1&0&0\\0&2&0&0\\0&0&2&1\\0&0&0&2}$.
		
		\begin{enumerate}[label=(\alph*)]
			\item Identify the eigenvalues of $C$ and their algebraic and geometric multiplicities. \hfill\textit{(2)}
			\item Write the Jordan canonical form $J$ of $C$. Draw the Jordan block structure, labelling each block with its eigenvalue and size. \hfill\textit{(3)}
			\item Let $M = C - 2I$. Show that $M$ is nilpotent and find its index of nilpotency. \hfill\textit{(2)}
			\item State one practical application of the Jordan form (e.g.\ in differential equations or matrix powers). \hfill\textit{(1)}
		\end{enumerate}
	\end{tcolorbox}
	
	\begin{tcolorbox}[answer]
	\end{tcolorbox}
	
	\newpage
	
	%=============================================================
	\section*{Formula Sheet}
	\addcontentsline{toc}{section}{Formula Sheet}
	%=============================================================
	
	\begin{tcolorbox}[colback=gray!5, colframe=gray!40, fonttitle=\bfseries, title={Key Formulae — for Reference}]
		
		\textbf{Determinants}
		\[
		\det\mat{a&b\\c&d} = ad-bc, \qquad
		\det(A) = \sum_{j=1}^{n}a_{ij}(-1)^{i+j}M_{ij} \text{ (cofactor expansion along row }i)
		\]
		
		\textbf{Cramer's Rule}\quad $x_i = \dfrac{\det(A_i)}{\det(A)}$, \;where $A_i$ replaces column $i$ of $A$ with $\mathbf{b}$.
		
		\textbf{Inverse of $2\times2$}\quad $\mat{a&b\\c&d}^{-1} = \dfrac{1}{ad-bc}\mat{d&-b\\-c&a}$
		
		\textbf{Orthogonal Projection}\quad
		$\text{proj}_W\mathbf{b} = \displaystyle\sum_{i=1}^{k}\frac{\mathbf{b}\cdot\mathbf{u}_i}{\mathbf{u}_i\cdot\mathbf{u}_i}\mathbf{u}_i$
		\quad (where $\{\mathbf{u}_1,\ldots,\mathbf{u}_k\}$ is an orthogonal basis for $W$)
		
		\textbf{Eigenvalues}\quad $\det(A - \lambda I) = 0$;\quad eigenvectors satisfy $(A-\lambda I)\mathbf{x}=\mathbf{0}$.
		
		\textbf{Trace–Eigenvalue relations}\quad $\text{tr}(A)=\sum\lambda_i$,\quad $\det(A)=\prod\lambda_i$
		
		\textbf{Diagonalisation}\quad $A = PDP^{-1}$,\quad $A^k = PD^kP^{-1}$
		
		\textbf{Matrix Exponential}\quad $e^A = \displaystyle\sum_{k=0}^{\infty}\frac{A^k}{k!}$;\quad if $A=PJP^{-1}$, then $e^A = Pe^JP^{-1}$
		
		\textbf{Rank-Nullity Theorem}\quad $\text{rank}(T)+\text{nullity}(T)=\dim(\text{domain})$
		
		\textbf{Matrix Norms}\quad $\norm{A}_\infty = \max_i\sum_j|a_{ij}|$,\quad $\norm{A}_1=\max_j\sum_i|a_{ij}|$
		
		\textbf{Nilpotent operator}\quad $N^k=O$ for some positive integer $k$ (index of nilpotency = smallest such $k$)
		
		\textbf{Jordan Block} of size $m$ for eigenvalue $\lambda$:
		\[
		J_m(\lambda) = \mat{\lambda & 1 & & \\ & \lambda & 1 & \\ & & \ddots & 1\\ & & & \lambda}_{m\times m}
		\]
		
	\end{tcolorbox}
	
\end{document}